\section{The congruence visibility filtration}\label{sec:filtration}

Let $\Sigma$ be a finite finitary signature with at least one constant,
and let $T$ be its closed term algebra. Constants have depth zero and
\[
\depth f(t_1,\ldots,t_k)=1+\max_i\depth t_i\qquad(k>0).
\]
Write $\equiv_E$ for full provable equality and
$\mathbb I_E=T/{\equiv_E}$ for the initial algebra.
Unless stated otherwise the language is one-sorted. The finite-ground
constructions apply sort by sort to a finite many-sorted signature; the
action example uses the two sorts of Paper I.

\begin{definition}[Bounded term layers and congruence proofs]
Put $T^{(D)}=\{t:\depth t\le D\}$. For ground $E$, let $\equiv_D$ be the
least equivalence relation on $T^{(D)}$ containing every input equation
whose two sides have depth at most $D$ and satisfying
\[
s_i\equiv_D t_i\ (1\le i\le k)
\quad\Longrightarrow\quad
f(s_1,\ldots,s_k)\equiv_D f(t_1,\ldots,t_k)
\]
whenever both displayed result terms belong to $T^{(D)}$.
Thus $\equiv_D$ is congruence on the induced partial term algebra.
\end{definition}

Equivalently, all terms in an equality proof DAG must have depth at most
$D$, including its premises and conclusions. Congruence can reuse a proved
argument equality without placing every intermediate term of that proof
inside the new context. For a finite universal presentation, replace the
input equations by all their ground instances with sides of depth at most
$D$. Sequential rewrites receive a separate definition in
Section~\ref{sec:rewrite}.

\begin{definition}[Observed, internal, and horizon equality]
For $D\ge d$, define the labelled quotient sets
\begin{align*}
I_E^{[d]}&=T^{(d)}/({\equiv_E}|_{T^{(d)}}),\\
J_E^{(d)}&=T^{(d)}/\equiv_d,\\
K_E^{(d,D)}&=T^{(d)}/({\equiv_D}|_{T^{(d)}}).
\end{align*}
They are not asserted to be total subalgebras. There are canonical
surjections $J_E^{(d)}\twoheadrightarrow K_E^{(d,D)}
\twoheadrightarrow I_E^{[d]}$.
\end{definition}

Whenever we write $K_E^{(d,D)}\cong I_E^{[d]}$, we mean that this canonical
surjection is bijective, or equivalently that the two equivalence
relations on the observed terms agree.

\begin{proposition}[Eventual visibility]\label{prop:eventual}
For every fixed $d$, the canonical map
$K_E^{(d,D)}\to I_E^{[d]}$ is bijective for all sufficiently large $D$.
\end{proposition}
\begin{proof}
There are finitely many observed pairs. Every full equality among them
has a finite proof containing finitely many terms. Take the maximum of
their depths and $d$. All these proofs are available by that horizon.
\end{proof}

\begin{definition}[Layer and pairwise visibility]
The exact horizon and its overhead are
\[
\delta_E(d)=\min\{D\ge d:{\equiv_D}|_{T^{(d)}}
 ={\equiv_E}|_{T^{(d)}}\},\qquad
\omega_E(d)=\delta_E(d)-d.
\]
For $d(s,t)=\max\{\depth s,\depth t\}$, put
\[
\lambda_E(s,t)=\min\{D\ge d(s,t):s\equiv_D t\},
\]
with $\lambda_E(s,t)=\infty$ if $s\not\equiv_E t$.
In particular $\lambda_E(t,t)=\depth t$.
\end{definition}

The layer invariant is the finite maximum
\[
\delta_E(d)=\max\left(d,
\max_{\substack{s,t\in T^{(d)}\\s\equiv_E t}}\lambda_E(s,t)\right).
\]
These are invariants of a \emph{presentation with a fixed proof convention}.
Adding a derived shallow equation can change them without changing
$\mathbb I_E$ or its equality relation on the original terms.

\subsection{Finite observation sets and coalescence levels}

For a finite query set $\mathcal Q$, set
$d_{\mathcal Q}=\max_{q\in\mathcal Q}\depth q$, taking $d_\varnothing=0$,
and define
\[
\delta_E(\mathcal Q)=\min\{D\ge d_{\mathcal Q}:
{\equiv_D}|_{\mathcal Q}={\equiv_E}|_{\mathcal Q}\}.
\]
Thus $\delta_E(\varnothing)=0$. In general, $\delta_E(d)=\max\{d,\delta_E(T^{(d)})\}$; when
$T^{(d)}$ contains a term of depth exactly $d$, it is simply the query-set special case.

\begin{proposition}[First-coalescence inequality]\label{prop:ultra}
If $s,t,u$ belong to one full equality class, then
\[
\lambda_E(s,u)\le
\max\{\lambda_E(s,t),\lambda_E(t,u)\}.
\]
\end{proposition}
\begin{proof}
Both premise equalities hold at the horizon on the right. Transitivity
at that horizon gives the conclusion.
\end{proof}

The nested partitions therefore have a merge-tree description. This
resembles an ultrametric hierarchy, but $\lambda_E(t,t)=\depth t$ need
not be zero, so $\lambda_E$ is not an ordinary ultrametric. For a final
query block $B$ and any $b\in B$, its coalescence level is
\[
\max_{q\in B}\lambda_E(b,q)
=\max_{p,q\in B}\lambda_E(p,q).
\]
The equality follows by the proposition, with endpoint activation included.

\begin{center}
\begin{tabularx}{\textwidth}{@{}lX@{}}
\toprule
Notation & Meaning \\
\midrule
$d$, $D$, $h(E)$ & Observation depth, proof horizon, maximum input-side depth.\\
$\lambda_E$, $\delta_E$ & Congruence-proof visibility; the finite-ground algorithm computes these.\\
$\lambda_E^{\mathrm{rw}}$, $\delta_E^{\mathrm{rw}}$ & Sequential rewrite visibility, defined in Section~\ref{sec:rewrite}.\\
$\mathcal Q$, $S$, $S_D$ & Query terms, their fixed input/subterm DAG, and its active nodes.\\
$\approx_D^S$ & Congruence closure on the active partial DAG.\\
$\mathsf P_D(E,\mathcal Q)$ & The partial quotient $S_D/\approx_D^S$.\\
$Q_E$, $\mathcal I$ & The carrier quotient and named action interface of Paper I.\\
\bottomrule
\end{tabularx}
\end{center}
